\newpage
\section{States and Observables}

\begin{abstract}

\noindent We begin by introducing the \emph{Data of States and Observables} (DOSO) framework, a foundational axiomatic system for the mathematical description of physical systems. DOSO aims to unify core concepts of states, observables, measurement, and dynamics across different physical theories. It axiomatizes a state space $\mcS$ as a real convex set (pure states being its extreme points) and a space of potential observables $\mcO$ as a complex topological vector space with structures like a real structure, a Lie algebra on a dense subspace, and a functional calculus. A measurement pairing links observables and states to probability distributions, and a mechanism for dynamics describes transformations generated by observables. \\

\noindent We then illustrate how to apply the DOSO framework to classical mechanics. Here, pure states correspond to points in phase space, observables to real-valued functions, the Lie algebra to Poisson brackets, and dynamics to Hamiltonian flows. The framework naturally accommodates statistical classical mechanics through probability measures on phase space.\\

\noindent Finally, we discuss how to apply DOSO to quantum mechanics. States are represented by density operators on a Hilbert space (pure states by rank-one projectors), observables by self-adjoint operators, the Lie algebra by commutators, and dynamics by unitary evolution. The abstract DOSO structure provides a common language compare classical and quantum mechanics. \\

\end{abstract}

\setcounter{tocdepth}{3}
\localtableofcontents

\textbf{References}: 
For foundational texts on the mathematical foundations of quantum mechanics, see \cite{Dirac1981}, \cite{vonNeumann1955MFQM}, \cite{Segal1947}, \cite{Mackey1963MFQM}. For more modern treatments of quantum theory, see \cite{Hall2013Quantum}, \cite{Strocchi2005}, and \cite{FaddeevYakubovskii2009}. For essential mathematical tools like functional analysis, see \cite{Conway1990}, \cite{ReedSimon1980}, \cite{Rudin1991}, \cite{vonNeumann1961}, \cite{Davidson1996}, and \cite{Murphy1990}. For classical mechanics, including Hamiltonian systems and symplectic geometry, see \cite{AbrahamMarsden1978}, \cite{CannasDaSilva2001}, and \cite{McDuffSalamon2017}. For geometric tools related to Lie groups, see \cite{Warner1983} and \cite{Hall2015Lie}.

\newpage 

\subsection{Our Basic Axioms}

\subsubsection{Data of States and Observables (DOSO)}

Describing physical reality mathematically requires a precise language, which we develop here. Central to any physical theory are the notions of \textit{state} (a complete description of a system at a given time), \textit{observable} (a measurable property), and \textit{dynamics} (how states or observables evolve in time). We adopt a general axiomatic framework encompassing classical mechanics, statistical mechanics, and, of particular interest, quantum mechanics and quantum field theory.

A physical system consists of its states, observables, and the pairing that assigns a value to each observable in a given state. For example, fix an integer $k > 0$ and consider $k$ particles moving in Euclidean space. A state is the collection of the particles' positions and velocities at a fixed time. An observable is a real-valued function on trajectories: for instance, the distance between the 2nd and 13th particles at 4:45 AM on January 10, 2023.

Now that we have the intuition down, we describe our axiomatization of states and observables.

\begin{definition}[Data of States and Observables (DOSO)] \label{ax:DOSO}
The data describing a physical system in terms of its states and observables comprises the following structures and their interactions:

\begin{enumerate}
\item (States.) A real convex space $\mcS$. Elements of $\mcS$ are \emph{states}, extreme points of $\mcS$ are \emph{pure states}, and non-extreme points are \emph{mixed states}. Let $\PS \subset \mcS$ denote the subset of pure states.

\item (Observables.) A complex topological vector space $\mcO$ called the \emph{algebra of observables}\footnote{It is $\cO_\R$ whose elements are observables. We consider the larger vector space $\cO$ for mathematical convenience.}, equipped with a \emph{real structure}, i.e., an antilinear involution $A\mapsto A^{*}$. Elements of
\[\mcO_{\R} = \{A \in \mcO \mid A^* = A\},\]
the subspace of real (or self-adjoint) elements of $\mcO$, are called \emph{observables}. There are the following additional structures.
\begin{enumerate}
    \item A \emph{measurement pairing}:
\[ \mcO_{\R}\times\mcS\longrightarrow \text{Prob}(\R) \]
\[ (A, \sigma)\mapsto\sigma_{A} \quad \]
where $\text{Prob}(\R)$ is the space of probability measures on $\R$ (equipped with its Borel $\sigma$-algebra).

    \item A map, called \emph{functional calculus} for observables:
\[ \text{Borel}(\R;\R)\times\mcO_{\R}\longrightarrow\mcO_{\R} \]
\[ (f,A)\mapsto f(A). \quad \]
\end{enumerate}

\item (Dynamics.) A dense subspace $\mcO^{\infty} \subseteq \mathcal{O}$ equipped with a complex Lie algebra structure. Let $\mcO_{\R}^{\infty}$ be the space of real elements of $\mcO^\infty$. For each $A \in \mcO_{\R}^{\infty}$, there exist
\begin{enumerate}
    \item a one-parameter group of automorphisms of $\mcS$;
    \item a one-parameter group of automorphisms of $\mcO$.
\end{enumerate}
\end{enumerate}
\end{definition}

Now, we briefly explain the purpose of each part of the axiomatization. Simultaneously, we introduce the full set of mathematical tools that this axiom provides us access to.

\subsubsection{Axiom 1: The Space of States ($\mcS$)}

\textit{Axiom 1: A real convex space $\mcS$. Elements of $\mcS$ are states, extreme points of $\mcS$ are pure states, and non-extreme points are mixed states. Let $\PS \subset \mcS$ denote the subset of pure states.}

\begin{remark}
Typically, $\mcS$ is a convex subset of a real topological vector space.
\end{remark}

\begin{remark}
In addition to making possible a unified treatment of pure states and mixed states, the introduction of a convex structure opens the door to tools from convexity theory, such as the Krein-Milman Theorem\footnote{A compact convex subset of a Hausdorff locally convex topological vector space is equal to the closed convex hull of its extreme points.} and Choquet's Theorem\footnote{Let $C$ be a compact convex subset of a normed space $V$. Given $c \in C$, there exists a probability measure $w$ supported on the set $E$ of extreme points of $C$ such that, for any affine function $f$ on $C$,
\[
f(c) = \int f(e) \, dw(e).
\]}.
\end{remark}

Before we discuss why convex structure constitutes part of the \dq{correct axiomatization,} let's first define all of the relevant terms.

\begin{definition}
\,
\begin{itemize}
\item Let $V$ be a real vector space. A subset $C \subseteq V$ is called a \textit{convex set} if for any two elements $x, y \in C$ and for any real number $t \in [0, 1]$, the element $tx + (1-t)y$ is also in $C$. The expression $t x + (1-t)y$ is called a \textit{convex combination} of $x$ and $y$.
\item  Let $C$ be a convex set. An element $e \in C$ is called an \textit{extreme point} of $C$ if it cannot be written as a non-trivial convex combination of two \textit{distinct} points in $C$. More formally, $e \in C$ is an extreme point if for any $x, y \in C$ and any $t \in (0, 1)$ (note the open interval), the condition $e = tx + (1-t)y$ implies that $x = y = e$. A \textit{non-extreme point} of a convex set $\mcS$ is any element $\sigma \in \mcS$ that is not an extreme point.
\end{itemize}
\end{definition}

%The concept of a real convex space naturally arises in quantum mechanics and statistical mechanics. Although physical states are often described using complex numbers (ie. quantum wavefunctions or density matrices), the outcomes of measurements, such as probabilities and expectation values, are always real numbers. Consequently, the underlying space of states is typically parameterized over the real numbers. For instance, while a density matrix in quantum mechanics may contain complex entries, it is Hermitian, and the convex combinations that represent statistical mixtures of states are formed using real coefficients.

%Convexity provides the mathematical framework that enables the formation of statistical mixtures. If $\sigma_1, \sigma_2 \in \mcS$ are two states, and $p \in [0,1]$ is a real number, then $p\sigma_1 + (1-p)\sigma_2$ represents a new state where the system is in state $\sigma_1$ with probability $p$ and in state $\sigma_2$ with probability $1-p$. This is very important in statistical mechanics and for describing systems where our knowledge is incomplete or for ensembles of systems.


Given two states $\sigma_1$ and $\sigma_2$ with probabilities $p_1$ and $p_2$ such that $p_1 + p_2 = 1$, there should be a state $\sigma'$ representing the system being in $\sigma_1$ with probability $p_1$ and in $\sigma_2$ with probability $p_2$. We call this a statistical mixture and write it as $p_1 \sigma_1 + p_2 \sigma_2$. Although it may not be obvious at first that statistical mixtures can be treated as linear combinations of states, this approach works well due to key properties of these mixtures. For instance, $p_1 \sigma_1 + p_2 \sigma_2 = p_2 \sigma_2 + p_1 \sigma_1$. Statistical mixtures play a fundamental role in statistical physics, and the framework of convex sets provides the proper mathematical setting for understanding and manipulating these combinations.

% because mixing $\sigma_1$ with probability $p_1$ and $\sigma_2$ with probability $p_2$ is naturally symmetric

%\rauan{Given two states $\sigma_1$ and $\sigma_2$ and probabilities $p_1, p_2$ with $p_1 + p_2 = 1$, there should be a state $\sigma'$ where the system is in state $\sigma_1$ with probability $p_1$ and in state $\sigma_2$ with probability $p_2$. We write the statistical mixture as $p_1 \, \sigma_1 + p_2 \, \sigma_2$. While it is not a priori clear that thinking of statistical mixtures as linear combinations of states is going to be successful, it does turn out to be successful thanks to some nice properties of statistical mixture. For example, $p_1 \, \sigma_1 + p_2 \, \sigma_2 = p_2 \, \sigma_2 + p_1 \, \sigma_1$ because being in state $\sigma_1$ with probability $p_1$ and being in state $\sigma_2$ with probability $p_2$ is the same as being in the state $\sigma_2$ with probability $p_2$ and being in state $\sigma_1$ with probability $p_1$. Statistical mixtures are important in statistical physics, and convex structure provides the right framework for $p_1 \, \sigma_1 + p_2 \, \sigma_2$ to make sense.}

% \textcolor{red}{The concept of a real space arises naturally in both quantum and statistical mechanics. While individual physical states (such as quantum wavefunctions or density matrices) may involve complex numbers, observable quantities like probabilities and expectation values are always real. Consequently, the space of physical states is typically parameterized over the real numbers.}

% \textcolor{red}{A key feature of this space is its convex structure, which allows for the formation of statistical mixtures. If $\sigma_1$ and $\sigma_2$ are possible states of a system and $p \in [0,1]$, then the combination}
% \[
% p \sigma_1 + (1-p)\sigma_2
% \]
% \textcolor{red}{represents a state where the system is in $\sigma_1$ with probability $p$ and in $\sigma_2$ with probability $1 - p$. This operation is commutative and distributive, reflecting the behavior of probabilistic combinations. For example, if $\sigma' = p_1 \sigma_1 + p_2 \sigma_2$, then mixing it with another state $\sigma_3$ using weights $q_1, q_2$ gives:}
% \[
% q_1 \sigma' + q_2 \sigma_3 = q_1 p_1 \sigma_1 + q_1 p_2 \sigma_2 + q_2 \sigma_3.
% \]
% \textcolor{red}{This convex structure forms the mathematical foundation for describing ensembles or systems with incomplete information.
% }





% \rauan{The concept of a real convex space arises naturally from probabilistic considerations. If $\sigma_1$ and $\sigma_2$ are two possible states of a physical system, and if $0 \le p_1, p_2 \le 1$ are real numbers with $p_1 + p_2 = 1$, there should be a notion of the system being in state $\sigma_1$ with probability $p_1$ and in $\sigma_2$ with probability $p_2$. Let us informally write this new state as
% \begin{equation*}
%     \sigma' = \text{($\sigma_1$ with $p_1$ and $\sigma_2$ with $p_2$)}.
% \end{equation*}
% Clearly, such probabilistic creation of states is commutative,
% \begin{equation*}
%     \text{($\sigma_1$ with $p_1$ and $\sigma_2$ with $p_2$)} = \text{($\sigma_2$ with $p_2$ and $\sigma_1$ with $p_1$)}.
% \end{equation*}
% There is also a kind of distributivity we have. Suppose $p_1 + p_2 = 1$ and $q_1 + q_2 = 1$, and let $\sigma'$ represent the system being in state $\sigma_1$ with probability $p_1$ and $\sigma_2$ with probability $p_2$. If $\sigma_3$ is any other state, then
% \begin{equation*}
%     \text{($\sigma'$ with $q_1$ and $\sigma_3$ with $q_2$)} = \text{($\sigma_1$ with $q_1 \, p_1$ and $\sigma_2$ with $q_1 \, p_2$ and $\sigma_3$ with $q_2$)}.
% \end{equation*}
% These are precisely the properties we would expect from a linear combination $p_1 \, \sigma_1 + p_2 \, \sigma_2$! Thus, we want to write the probabilistic combination as
% \begin{equation}
%     \boxed{\text{($\sigma_1$ with $p_1$ and $\sigma_2$ with $p_2$)} = p_1 \, \sigma_1 + p_2 \, \sigma_2.}
% \end{equation}
% A real convex space is precisely the mathematical framework we need to make this linear combination make sense. Convex structure allows us to take \emph{statistical mixtures} of states, which is of great importance in statistical physics.}

We choose the extreme points of the convex space to be our pure states because physically, pure states represent states of maximal knowledge. They are \dq{indecomposable} in the sense that they cannot be viewed as statistical mixtures of other, different states. Identifying pure states is important as they often form the \dq{building blocks} of the theory. Then, it naturally follows that we should define the mixed states as the non-extreme points.

\begin{example}[Classical mechanics]
In classical statistical mechanics, states are probability measures on the phase space, and the space of such measures is convex. A pure state is a Dirac measure at some phase space point, representing the system certainly being in that phase space configuration.
\end{example}

\begin{example}[Quantum mechanics]
In quantum mechanics, states are density operators $\rho$ on a Hilbert space $\mcH$. These are positive trace-class operators\footnote{Let $\mathcal{H}$ be a separable Hilbert space. A (not necessarily positive) bounded operator $T:\mathcal{H}\to \mathcal{H}$ is called \textit{trace-class} if and only if $\text{Tr}(|T|) <\infty$ where $|T| :=\sqrt{T^* T}$.} with $\Tr(\rho)=1$. The set of such operators is convex. A pure state is equal to a projection $\ketbra{\psi}{\psi}$ for some unit vector $\ket{\psi} \in \cH$. Such a state cannot be written as $\rho = p \, \rho_1 + (1-p) \, \rho_2$ for $0 < p < 1$ and $\rho_1, \rho_2 \neq \rho$ being other density operators.
\end{example}

Let's move onto the next axiom.

\subsubsection{Axiom 2: The Algebra of Observables ($\mcO$)}
\textit{Axiom 2: A complex topological vector space $\mcO$ called the \emph{algebra of observables}\footnote{It is $\cO_\R$ whose elements are observables. We consider the larger vector space $\cO$ for mathematical convenience.}, equipped with a \emph{real structure}, i.e., an antilinear involution $A\mapsto A^{*}$. Elements of
\[\mcO_{\R} = \{A \in \mcO \mid A^* = A\},\]
the subspace of real (or self-adjoint) elements of $\mcO$, are called \emph{observables}. There are the following additional structures.
\begin{enumerate}
    \item A \emph{measurement pairing}:
\[ \mcO_{\R}\times\mcS\longrightarrow \textnormal{Prob}(\R) \]
\[ (A, \sigma)\mapsto\sigma_{A} \quad \]
where $\textnormal{Prob}(\R)$ is the space of probability measures on $\R$ (equipped with its Borel $\sigma$-algebra).
    \item A map, called \emph{functional calculus} for observables:
\[ \textnormal{Borel}(\R;\R)\times\mcO_{\R}\longrightarrow\mcO_{\R} \]
\[ (f,A)\mapsto f(A). \quad \]
\end{enumerate}
}

The space $\mcO$ and its associated structures provide the mathematical representation of physical quantities that can be measured. There is quite a few properties to discuss, and we survey them one by one.

%The choice of $\mcO$ as a complex topological vector space is motivated by physics, with some additional mathematical structure added because it gives us access to more tools. Although measurement outcomes are real, quantum mechanics is naturally framed in complex Hilbert spaces, where observables are self-adjoint operators. Complex numbers offer a richer algebraic structure essential for tools like spectral theory and Fourier transforms. In classical systems, complexifying function spaces, such as $C(X, \C)$, simplifies analysis. The requirement for $\mcO$ to be a topological vector space arises from the need for continuity, convergence, and approximation in physics. Many observables, like position and momentum, are unbounded operators, requiring operator topologies (norm, strong, weak) for rigorous treatment. This topological structure enables the use of functional analysis, the most useful tool for infinite-dimensional systems in quantum mechanics and field theory.

\begin{definition}[Complex Topological Vector Space]
A \textit{complex topological vector space} is a vector space over the field of complex numbers $\C$ that is endowed with a topology such that vector addition $(u,v) \mapsto u+v$ and scalar multiplication $(\lambda, v) \mapsto \lambda v$ are continuous functions.
\end{definition}

\textbf{Topological vector space structure.} If $A_1$ and $A_2$ are observables and $\lambda$ a real number, then, of course, we can define observables $A_1 + A_2$ and $\lambda \, A_1$ -- therefore, the mathematical object containing all observables should be a real vector space. Thus, we put observables into the real vector space $\cO$.

Approximation is an invaluable technique in physics. Therefore, we require $\cO$ to be a \emph{topological} vector space, meaning that there is a reasonable notion of continuity and limits. Any finite-dimensional vector space has an obvious topology, namely, the Euclidean one -- yet one has to be careful in the infinite-dimensional case. Many observables, like position and momentum, are unbounded operators, requiring operator topologies (norm, strong, weak) for rigorous treatment. This topological structure enables the use of functional analysis, the most useful tool for infinite-dimensional systems in quantum mechanics and field theory.

\textbf{Complex vector space structure.} We go a step further and say that $\cO$ is a \emph{complex} vector space.  Making $\mcO$ a \emph{complex} vector space is motivated partly by quantum mechanics, a theory that is naturally expressed in complex vector spaces. Considering complex numbers is also mathematically convenient. Firstly, they offer a richer algebraic structure essential for tools like spectral theory and Fourier transforms. Secondly, complexifying function spaces, such as $C(X, \C)$, simplifies analysis.

\begin{example}
In classical mechanics, $\mcO$ might be a space of functions on phase space, like $C(M)$ (continuous functions on the phase space manifold $M$) or $L^p(M, d\mu)$ (functions whose $p$-th power is integrable with respect to a measure $\mu$), which are often Banach spaces and thus topological vector spaces. In quantum mechanics, $\mcO$ is typically related to the algebra $B(\mathcal{H})$ of bounded linear operators on a Hilbert space $\mathcal{H}$, or more generally, a $C^*$-algebra or a von Neumann algebra.
\end{example}

\begin{definition}[Real Structure on a Complex Vector Space]
A \textit{real structure} on a complex vector space $V$ is an antilinear map $J: V \to V$ (often denoted by $A \mapsto A^*$) such that $J^2 = \text{id}$ (i.e., $J(J(A)) = (A^*)^* = A$ for all $A \in V$). An antilinear map $J$ satisfies $J(\lambda A + \mu B) = \bar{\lambda} J(A) + \bar{\mu} J(B)$ for all $A, B \in V$ and $\lambda, \mu \in \C$. The set of fixed points of $J$, $V_{\R} = \{A \in V \mid J(A) = A\}$, forms a real vector space.
\end{definition}

\textbf{Real structure.} Complexifying $\cO$ comes with the obvious drawback that not every $A \in \cO$ is an observable. Intuitively, the complex structure of $\cO$ would conflict with the fact that outcomes of physical measurements are always real numbers. (We will discuss measurement in later.) Complex numbers are nice to work with, yet the outcomes of physical measurements are always real numbers. The usefulness of a real structure lies in its ability to bridge the potentially complex algebraic framework of $\mcO$ with the real-valued nature of physical measurements.

\begin{example}
In classical mechanics, if $\mcO$ is taken as the space of complex-valued functions on phase space $M$, $C(M, \C)$, then the star operation $A^*$ corresponds to pointwise complex conjugation $f \mapsto \bar{f}$, and observables $\mcOR$ are the real-valued functions $C(M, \R)$. In quantum mechanics, if $\mcO$ is an algebra of operators on a Hilbert space $\mathcal{H}$ (like $B(\mathcal{H})$), then $A^*$ is the standard operator adjoint, and $\mcOR$ comprises the self-adjoint (Hermitian) operators.
\end{example}

This axiom formalizes the important connection between the abstract mathematical objects of the theory (states and observables) and the empirical results obtained from physical measurements. It essentially defines how to extract predictions for experimental outcomes from the theory.

\begin{definition}[Probability Measure on $\R$]
A \textit{probability measure on $\R$} (more precisely, on the Borel $\sigma$-algebra of $\R$, denoted $\text{Borel}(\R)$) is a function $P: \text{Borel}(\R) \to [0,1]$ such that:
\begin{enumerate}
    \item $P(E) \ge 0$ for all $E \in \text{Borel}(\R)$ (non-negativity).
    \item $P(\R) = 1$ (normalization).
    \item For any countable sequence of pairwise disjoint Borel sets $E_1, E_2, \dots$, $P(\bigcup_{i=1}^\infty E_i) = \sum_{i=1}^\infty P(E_i)$ (countable additivity or $\sigma$-additivity).
\end{enumerate}
The space of all such probability measures on $\R$ is denoted by $\text{Prob}(\R)$.
\end{definition}

\textbf{Measurement pairing.} This bullet point should come rather naturally. Observables are physical quantities that can be measured. If some observable is measured in some state, the outcome is not always deterministic -- in quantum mechanics, for example, measurement includes an intrinsic element of randomness. For this reason, we say that measuring an observable $A$ in a state $\sigma$ produces results distributed according to some distribution $\sigma_A$.

%The measurement pairing takes a real observable $A \in \mcOR$ and a state $\sigma \in \mcS$ as input, and outputs a probability measure $\sigma_A$ on the real line. The domain $\mcOR \times \mcS$ means that predictions are made for specific observables when the system is in a specific state. The choice of $\mcOR$ for the observable is because (as established previously) physical measurements are expected to yield real-numbered outcomes, and $\mcOR$ represents precisely those observables (the self-adjoint elements) corresponding to such quantities. The codomain $\text{Prob}(\R)$ is motivated by physics: the outcome of measuring an observable $A$ for a system in state $\sigma$ is generally not deterministic. Instead, the theory predicts a probability distribution $\sigma_A$ over the possible real-valued outcomes for $A$. This distribution $\sigma_A$ encapsulates all statistical information about the measurement of $A$ in state $\sigma$; for instance, $\sigma_A(E)$ is the probability that the measurement outcome lies within the Borel set $E \subset \R$.

\begin{example}
In classical mechanics, if the state $\sigma$ is a pure state, corresponding to a specific point $x_0$ in phase space, and $A$ is an observable represented by a function $A(x)$ on phase space, the measurement is deterministic. The resulting probability measure $\sigma_A$ is then the Dirac delta measure $\delta_{A(x_0)}$, which assigns probability 1 to the single value $A(x_0)$ and 0 to any set not containing $A(x_0)$. If $\sigma$ is a classical mixed state, represented by a probability measure $\mu$ on phase space (e.g., a Liouville distribution), then $\sigma_A$ is the pushforward measure $A_*\mu$, defined by $(A_*\mu)(E) = \mu(A^{-1}(E))$ for any Borel set $E \subset \R$. This gives the probability distribution for the values of $A(x)$ when $x$ is sampled according to $\mu$. In quantum mechanics, the probabilistic nature is even more apparent. If $\sigma$ is a pure state represented by a normalized vector $|\psi\rangle$ in a Hilbert space $\mathcal{H}$, and $A$ is a self-adjoint operator, the spectral theorem associates $A$ with a unique projection-valued measure (PVM) $\Pi_A$. The probability that a measurement of $A$ yields a value in a Borel set $E \subset \R$ is then given by $\langle\psi, \Pi_A(E)\psi\rangle$. This expression defines the measure $\sigma_A(E)$. For a quantum mixed state described by a density operator $\rho$, this generalizes to $\sigma_A(E) = \Tr(\rho \Pi_A(E))$. This formulation, often referred to as the Born rule (or its generalization), is very important in quantum theory.
\end{example}

%This axiom is very useful: It provides the explicit rule for calculating the probabilities of different measurement outcomes for any given observable in any given state. These calculated probabilities (or derived quantities like expectation values $\langle A \rangle_{\sigma} = \int x \, d\sigma_A(x)$ and variances) are what can be directly compared with experimental data gathered from repeated measurements on ensembles of identically prepared systems. Thus, this axiom bridges the gap between the mathematical formalism and empirical verification, allowing the theory to be tested, validated, or refuted.

\begin{definition}[Borel Measurable Function]
Let $(X, \mathcal{T})$ be a topological space, and let $\mathcal{B}(X)$ be the Borel $\sigma$-algebra on $X$ (the smallest $\sigma$-algebra containing all open sets in $\mathcal{T}$). Let $(Y, \mathcal{F})$ be a measurable space (a set $Y$ with a $\sigma$-algebra $\mathcal{F}$ of subsets of $Y$). A function $h: X \to Y$ is \textit{Borel measurable} (or simply Borel) if for every set $S \in \mathcal{F}$, the preimage $h^{-1}(S) = \{x \in X \mid h(x) \in S\}$ is in $\mathcal{B}(X)$.
For the functional calculus, $f: \R \to \R$, both the domain and codomain are equipped with their standard Borel $\sigma$-algebras.
\end{definition}

\textbf{Functional calculus.} Intuitively, if $A$ is an observable and $f$ is a function, it is natural that we can define an observable $f(A)$ as follows: for a physical measurement in which $A$ would have been measured to have value $a \in \R$, the observable $f(A)$ is $f(a)$. The definition of a functional calculus for observables only makes this intuition precise.%Restricting to Borel functions is a technical mathematician's condition to avoid unreasonable behaviors.

The functional calculus, a map $(f,A) \mapsto f(A)$ from $\text{Borel}(\R;\R)\times\mcOR$ to $\mcOR$, allows for the construction of new real observables from existing ones by applying real-valued Borel measurable functions. If $A$ is an observable representing a physical quantity, and $f: \R \to \R$ is such a function, then $f(A)$ is also considered a well-defined observable. The choice of Borel functions is significant because this class is sufficiently broad to include most functions of practical interest in physics, such as polynomials (e.g., forming $A^2$), exponentials (e.g., $e^{cA}$), and characteristic functions of intervals or other Borel sets (e.g., $\chi_E(A)$, which acts as a projector onto measurement outcomes of $A$ lying in the set $E$). Furthermore, the set of Borel functions is well-behaved with respect to pointwise limits of sequences of functions.

\begin{example}
In classical mechanics, if an observable $A$ is represented by a real-valued function $A(q,p)$ on phase space, then $f(A)$ is simply the composite function $(f \circ A)(q,p)$. In quantum mechanics, for a self-adjoint operator $A$, the functional calculus $f(A)$ is rigorously defined via the spectral theorem. If $A$ has a spectral decomposition $A = \int_{\sigma(A)} \lambda \, dP_\lambda$, where $P_\lambda$ is its spectral measure and $\sigma(A)$ its spectrum, then $f(A)$ is defined as $f(A) = \int_{\sigma(A)} f(\lambda) \, dP_\lambda$.
\end{example}


\subsubsection{Axiom 3: Dynamics and Symmetries}
\textit{Axiom 3: A dense subspace $\mcO^{\infty} \subseteq \mathcal{O}$ equipped with a complex Lie algebra structure. Let $\mcO_{\R}^{\infty}$ be the space of real elements of $\mcO^\infty$. For each $A \in \mcO_{\R}^{\infty}$, there exist
\begin{enumerate}
    \item a one-parameter group of automorphisms of $\mcS$;
    \item a one-parameter group of automorphisms of $\mcO$.
\end{enumerate}}

This axiom introduces the notion of evolution or transformation within the system, establishing how changes, such as time evolution or symmetry operations, are generated by certain well-behaved real observables. It provides the framework for dynamics and continuous symmetries.

\begin{definition}[Complex Lie Algebra]
A \textit{complex Lie algebra} is a vector space $L$ over the complex numbers $\C$, equipped with a binary operation $[\cdot, \cdot]: L \times L \to L$, called the Lie bracket, satisfying the following properties for all $A, B, C \in L$ and $\alpha, \beta \in \C$:
\begin{enumerate}
\item Bilinearity: $[\alpha A + \beta B, C] = \alpha[A, C] + \beta[B, C]$ and $[A, \alpha B + \beta C] = \alpha[A, B] + \beta[A, C]$.
\item Antisymmetry: $[A, B] = -[B, A]$.
\item Jacobi Identity: $[A, [B, C]] + [B, [C, A]] + [C, [A, B]] = 0$.
\end{enumerate}
\end{definition}

The existence of a complex Lie algebra structure is a natural axiom motivated by the evolution of observables in the Heisenberg picture, which we talk about later. In classical mechanics, the Poisson bracket $\{f, g\}$ plays a similar role in Hamiltonian dynamics and canonical transformations. The \dq{complex} nature of the Lie algebra makes the bracket $\C$-bilinear. Compatibility conditions like $[A_1, A_2]^* = [A_1^*, A_2^*]$ ensure the real structure interacts coherently with the Lie bracket, implying that the real subspace $\mcORinfty$ forms a real Lie subalgebra. This Lie algebra structure is essential for understanding the non-commutative nature of observables, central to both classical and quantum Hamiltonian mechanics, continuous symmetries (via Noether’s theorem), and quantization procedures that relate classical Poisson algebras to quantum Lie algebras.

\begin{definition}[Dense Subspace]
Let $X$ be a topological space and $S$ be a subspace of $X$ (i.e., $S \subseteq X$ with the induced topology). $S$ is a \textit{dense subspace} of $X$ if the closure of $S$ in $X$ is $X$ itself ($\bar{S} = X$). Equivalently, for every point $x \in X$ and every open neighborhood $U$ of $x$, $U \cap S \neq \emptyset$.
\end{definition}

Putting the Lie algebra structure a dense subspace $\mcOinfty \subset \mcO$ addresses the common situation in physics where certain algebraic operations (like forming Lie brackets) or analytical procedures (such as differentiation with respect to parameters) are well-defined or exhibit better behavior only on a smaller, more \dq{regular} or \dq{smoother} subset of observables. This $\mcOinfty$ represents such a well-behaved subset. The property of density is critical because it ensures that properties defined on $\mcOinfty$, or elements within it, can often be extended by continuity or approximation to the entirety of $\mcO$. Thus, this dense subspace provides a well-behaved domain for structures like the Lie algebra, which might not be definable or manageable on all of $\mcO$.

\begin{example}
In classical mechanics, if $\mcO$ is a space like $L^2(M)$ or $C_0(M)$ (continuous functions vanishing at infinity), $\mcOinfty$ might be $C^\infty(M)$, the space of smooth functions on phase space. Poisson brackets, important to Hamiltonian dynamics, are naturally and rigorously defined on $C^\infty(M)$, which is dense in many relevant larger function spaces. In quantum mechanics, if dealing with unbounded operators, $\mcOinfty$ could be a common, dense domain of important self-adjointness for a set of operators, or it might represent operators that are \dq{smooth} with respect to a given dynamical evolution. For instance, the commutator of two unbounded self-adjoint operators might not be densely defined or lead to another self-adjoint operator without careful consideration of domains, necessitating the selection of a suitable $\mcOinfty$.
\end{example}

\begin{definition}[One-Parameter Group of Automorphisms]
A \textit{one-parameter group of automorphisms} on a space $X$ (either $\mcS$ or $\mcO$) is a family of automorphisms $\{\phi_t\}_{t \in \R}$ indexed by a real parameter $t$ (often representing time or a continuous symmetry parameter) such that:
\begin{enumerate}
    \item $\phi_0 = \text{id}$ (the identity map on $X$).
    \item $\phi_t \circ \phi_s = \phi_{t+s}$ for all $t, s \in \R$ (group homomorphism property).
    \item The map $t \mapsto \phi_t(x)$ is continuous for each $x \in X$ (continuity property, with respect to the topology on $X$ and $\R$). The type of continuity (e.g., strong, weak) depends on the context and topology of $X$.
\end{enumerate}
\end{definition}

Physical transformations, such as time evolution, spatial translations, or rotations, are typically understood to be generated by specific real physical quantities: the Hamiltonian generates time evolution, momentum generates spatial translations, angular momentum generates rotations, and so on. The restriction of generators to $\mcORinfty$ ensures that these are indeed real observables and that they are \dq{sufficiently well-behaved} (due to being in $\mcOinfty$) for the generated transformations to be well-defined, for example, leading to differentiable paths for states or observables, or forming strongly continuous one-parameter groups as per Stone's theorem, which we discuss in Lecture 4.

By Axiom 3, every $A \in \mcOinfty$ has associated with it a one-parameter group of automotphisms of $\mcS$ (evolution of states) and another such group for $\mcO$ (evolution of observables). The first map, where a generator $A \in \mcORinfty$ gives rise to a one-parameter group of automorphisms $\{\alpha_t^A\}_{t\in\R}$ on $\mcS$, corresponds to the \textit{Schrödinger picture}. In this picture, the observables are considered fixed (or only explicitly time-dependent), while the state of the system $\sigma$ evolves according to $\sigma(t) = \alpha_t^A(\sigma(0))$. The second map, where $A$ generates a group $\{\beta_t^A\}_{t\in\R}$ on $\mcO$, corresponds to the \textit{Heisenberg picture}. Here, the state $\sigma$ is considered fixed, while the observables $A$ evolve according to $A(t) = \beta_t^A(A(0))$. For the framework to be consistent, these two pictures must yield equivalent physical predictions for the measurement outcomes $\sigma_A(t)$. This implies a compatibility condition between how states and observables transform, ensuring that the probability measure $\left(\sigma(t)\right)_A$ (Schrödinger) is the same as $\sigma_{\left(A(t)\right)}$ (Heisenberg).

\begin{example}
In quantum mechanics, time evolution is governed by the Hamiltonian operator $H \in \mcORinfty$. In the Schrödinger picture, state vectors evolve as $\ket{\psi(t)} = U_t \ket{\psi(0)}$, where $U_t = e^{-itH/\hbar}$ is a one-parameter unitary group. This induces automorphisms on the space of density operators (states $\mcS$) via $\rho(t) = U_t \rho(0) U_t^*$. Observables (operators) do not evolve in time unless explicitly specified. In the Heisenberg picture, observables evolve as $A(t) = U_t^* A(0) U_t$, while states do not evolve in time. Both $U_t(\cdot)U_t^*$ and $U_t^*(\cdot)U_t$ are one-parameter groups of automorphisms on the appropriate spaces.
\end{example}

\begin{example}
In classical mechanics, a Hamiltonian function $H \in C^\infty(M) \subset \mcORinfty$ generates a Hamiltonian flow $\phi_t$ on phase space $M$, which is a one-parameter group of symplectomorphisms, which we discuss in Lecture 2. This flow induces an evolution of states (probability measures on $M$) via the pushforward, $\mu_t = (\phi_t)_* \mu_0$, and an evolution of observables (functions on $M$) via the pullback, $f_t = f \circ \phi_{-t}$.
\end{example}

The usefulness of Axiom 3 lies in how it allows the theory to make predictions about how physical systems change over time. Beyond time evolution, it provides the mathematical machinery to describe continuous symmetries of the system. If a transformation leaves the system's properties invariant, its generator (an element of $\mcORinfty$) is often a conserved quantity associated with that symmetry. This is the concept behind Noether's theorem, linking symmetries to conservation laws, an important principle in physics.


\subsubsection{Summary of DOSO}
The following table summarizes the relationship between DOSO and three physical frameworks: classical mechanics, statistical mechanics, and quantum mechanics. The three physical theories are, of course, not mutually exclusive: statistical physics studies classical as well as quantum systems. We show how DOSO translates into the terms and conventions commonly used in classical mechanics, statistical mechanics, and quantum mechanics.

{\tiny
\begin{longtable}{|p{0.23\textwidth}|p{0.24\textwidth}|p{0.24\textwidth}|p{0.24\textwidth}|}
\hline
\textit{DOSO Component} & \textit{Classical Mechanics (CM)} & \textit{Statistical Mechanics (SM)} & \textit{Quantum Mechanics (QM)} \\
\hline
\endfirsthead
\multicolumn{4}{c}%
{{\bfseries Table \thetable\ continued from previous page}} \\
\hline
\textit{DOSO Component} & \textit{Classical Mechanics (CM)} & \textit{Statistical Mechanics (SM)} & \textit{Quantum Mechanics (QM)} \\
\hline
\endhead
\hline \multicolumn{4}{r}{{Continued on next page}} \\
\endfoot
\hline
\endlastfoot

\textit{1. Space of States ($\mcS$)} &
\textit{Pure:} Phase space point.
\newline
\textit{Mixed:} Probability measure on phase space. &
\textit{State:} Probability measure on phase space.
\newline
\textit{Pure (SM):} Dirac delta measure (microstate).
\newline
\textit{Mixed:} Ensemble. &
\textit{State:} Density operator.
\newline
\textit{Pure:} Rank-one projector.
\newline
\textit{Mixed:} Non-projector density operator.
\\
\hline

\textit{2. Space of Operators ($\mcO$) \& Structures} &
\textit{Observables:} Real functions on phase space.
\newline
\textit{Smooth Subspace:} Smooth functions.
\newline
\textit{Lie Algebra:} Poisson bracket.
\newline
\textit{Functional Calculus:} Function composition. &
\textit{Observables:} Real functions on phase space.
\newline
\textit{Smooth Subspace:} Smooth functions.
\newline
\textit{Lie Algebra:} Poisson bracket (for Liouville eq.).
\newline
\textit{Functional Calculus:} Function composition (for averages). &
\textit{Observables:} Self-adjoint operators.
\newline
\textit{Smooth Subspace:} E.g., Bounded operators or specific domains.
\newline
\textit{Lie Algebra:} Commutator.
\newline
\textit{Functional Calculus:} Defined via the spectral theorem.
\\
\hline

\textit{3. Measurement Pairing} &
Deterministic value for pure states. &
Pushforward measure for ensembles. &
Born Rule (probabilistic outcomes).
\\
\hline

\textit{4. Dynamics \& Symmetries} &
Hamilton's equations for phase space points.
\newline
Poisson bracket with the Hamiltonian for observables. &
Liouville's equation for probability density. &
Schrödinger equation for states (density operators/wavefunctions).
\newline
Heisenberg equation of motion for operators.
\\
\hline

\end{longtable}
}

\subsubsection{Compatibilities Among DOSO Data}
The structures outlined in Axiom System \ref{ax:DOSO} are not independent but must satisfy certain important compatibility conditions for the framework to be coherent and physically meaningful. The following are some examples of such requirements.

For $f\in \text{Borel}(\R;\R)$, $A, A_{1}, A_{2}\in\mcO$, $\sigma, \sigma_{i}\in\mcS,$ and positive real numbers $x^{i}$ (more accurately $p_i \in [0,1]$) with $\sum p_i = 1$ (for a convex combination $p_1\sigma_1 + p_2\sigma_2 + \dots$), we require:

\begin{enumerate}
    \item Linearity of measurement with respect to state mixing:
    If $\sigma = \sum_i p_i \sigma_i$ is a mixed state, then the probability distribution of $A$ in state $\sigma$ is the corresponding mixture of probability distributions:
    \[ \left(\sum_i p_i\sigma_{i}\right)_{A}=\sum_i p_i(\sigma_{i})_{A}\]
    
    This condition is important to the interpretation of mixed states as statistical ensembles. The probability distribution obtained from a mixed state is precisely the weighted average of the probability distributions obtained from its underlying states. This applies to the probability measures themselves, implying linearity for expectation values as well.
    
    \item Behavior of measurement for sums/functions of independent observables:
    \[ \sigma_{A_{1}+A_{2}}=\sigma_{A_{1}}*\sigma_{A_{2}}= \int_{-\infty}^{\infty} \sigma_{A_1}(y) \cdot \sigma_{A_2}(x - y) \, dy \]
    This states that the probability distribution for the sum of two independent observables $A_1, A_2$ is the convolution of their individual probability distributions. This is a standard probability assumption that is true in classical, statistical, and quantum mechanics. 
    
    \item Compatibility of the Lie bracket with the real structure:
    \[ [A_{1},A_{2}]^{*}=[A_{1}^{*},A_{2}^{*}] \quad \text{for } A_{1}, A_{2}\in\mcO^{\infty}\]
    
    This condition ensures that the real structure (the $*$-involution) is compatible with the Lie algebra structure on $\mcO^\infty$. It means that $*$ is an antilinear Lie algebra automorphism. An important consequence is that if $A_1, A_2$ are real observables (i.e., $A_1^*=A_1, A_2^*=A_2$), then $[A_1,A_2]^* = [A_1,A_2]$. This implies that $[A_1,A_2]$ is also a real observable, provided the Lie bracket itself maps pairs of real observables to real observables.
    
    In classical mechanics, the compatibility of the Poisson bracket with the real structure ensures that the evolution of observables in phase space preserves their real values, while maintaining the consistency of symmetries and conserved quantities (e.g., energy, momentum) under canonical transformations. In statistical mechanics, the real structure guarantees that macroscopic thermodynamic quantities, such as the partition function and free energy, remain real and well-defined, while ensuring that the statistical evolution respects the symmetries of the system (e.g., time-reversal or gauge invariance). In quantum mechanics, this ensures that commutation relations between Hermitian operators preserve the Hermitian property, thus guaranteeing real eigenvalues for physical observables.
    
\end{enumerate}

For $H\in\mcO_{\R}^{\infty}$ (a real, well-behaved observable, often the Hamiltonian or a symmetry generator), let
\begin{align*}
    \sigma&\mapsto\sigma^{(t)} \\
    A&\mapsto A^{(t)}
\end{align*}
denote the generated one-parameter groups of automorphisms of $\mcS$ and $\mcO$, respectively. Then:

\begin{enumerate}
    \item The Heisenberg equation of motion:
    \[ \frac{d}{dt}A^{(t)}+[H,A^{(t)}]=0, \quad \text{for } A\in\mcO^{\infty}, t\in\R \]
    
    This is the abstract form of the Heisenberg equation of motion. It describes how an observable $A$ (from the \dq{smooth} set $\mcO^\infty$) evolves in time, $A \to A^{(t)}$, under the influence of the generator $H$. The Lie bracket $[H, A^{(t)}]$ dictates the rate and manner of this change.
    
    The existence and uniqueness of solutions to this differential equation are very important, but we will not discuss it here.
    
    \item Equivalence of Schrödinger and Heisenberg pictures for measurement outcomes:
    \[ (\sigma^{(t)})_{A}=\sigma_{A^{(t)}}, \quad \text{for } \sigma\in \mcS, A\in\mcO, t\in\R \]
    
    This is a important consistency condition. It states that the probability distribution for measuring a fixed observable $A$ in an evolved state $\sigma^{(t)}$ (Schrödinger picture) is identical to the probability distribution for measuring the evolved observable $A^{(t)}$ in the initial state $\sigma$ (Heisenberg picture). This ensures that the physical predictions are independent of the \dq{picture} we artificially adopted, establishing a duality between the two pictures.
    
\end{enumerate}

\subsubsection{Expectation Values and Separability}

From the measurement pairing, we can define the \emph{expectation value} (or average measured value) of an observable.

\begin{definition}[Expectation Value]
Given a state $\sigma \in \mcS$ and a real observable $A \in \mcO_{\R}$, the probability measure $\sigma_A \in \text{Prob}(\R)$ describes the distribution of outcomes if $A$ is measured in state $\sigma$. The \emph{expectation value} of $A$ in state $\sigma$, denoted $\langle A \rangle_{\sigma}$, is defined as the mean of this distribution:
\[ \langle A \rangle_{\sigma} \coloneqq \int_{\R} \lambda \, d\sigma_{A}(\lambda) \]
provided the integral exists. This defines a real-valued pairing $\mcS\times\mcO_{\R}\rightarrow\R$.
\end{definition}


The expectation value $\langle A \rangle_{\sigma}$ represents the average result that would be obtained from a large number of independent measurements of the observable $A$ on identical copies of the system, each prepared in the state $\sigma$. It is one of the most important quantities extracted from the probability distribution $\sigma_A$.

An important property of this pairing is that it \dq{separates points}:
\begin{enumerate}
    \item If $\sigma_{1}, \sigma_{2}\in\mcS$ satisfy $\langle A\rangle_{\sigma_{1}}=\langle A\rangle_{\sigma_{2}}$ for all $A\in\mcO_{\R}$ for which the expectation values are defined, then $\sigma_{1}=\sigma_{2}$.
    
    This means that distinct states must be distinguishable by at least one observable's expectation value. If two states yield the same average outcome for every possible measurement, they are the same state. This ensures that the space of states $\mcS$ does not contain redundant information.
    
    \item If $A_{1}, A_{2}\in\mcO_{\R}$ satisfy $\langle A_{1}\rangle_{\sigma}=\langle A_{2}\rangle_{\sigma}$ for all $\sigma\in\mcS$ for which the expectation values are defined, then $A_{1}=A_{2}$ (in $\mcO_\R$).
    
    This is the same thing, but for observables. Distinct observables must be distinguishable by their expectation value in at least one state. If two observables yield the same average outcome in every possible state of the system, they are the same observable. This ensures that $\mcO_\R$ does not contain operationally redundant elements.
\end{enumerate}
This separation property implies that the pairing $\langle \cdot, \cdot \rangle : \mcS \times \mcO_\R \to \R$ is non-degenerate.

\subsubsection{Further Remarks on the Axiom System}

We make several remarks about Axiom System \ref{ax:DOSO}.

\begin{remark}[Interpretations and Implications] \label{rem:1.9}
\mbox{} % Empty box for correct spacing after remark title
\begin{enumerate}
\item Schrödinger vs. Heisenberg Pictures: The equation ${\sigma^{(t)}}_{A}=\sigma_{A^{(t)}}$ is the mathematical expression of the physical equivalence between the Schrödinger picture (where states evolve in time and observables are fixed) and the Heisenberg picture (where states are fixed and observables evolve). The laws of physics should not depend on which calculational framework one chooses to describe the evolution of measurement outcome probabilities.

\item Probabilistic Nature of Measurement: Measurement is inherently probabilistic in quantum mechanics. In classical mechanics, if a system is in a pure state $\sigma$ (e.g., a specific point in phase space detailing all positions and momenta), then the outcome $\sigma_A$ of measuring an observable $A$ (a function on phase space) is a point measure (Dirac delta measure at the value $A(\sigma)$). Thus, for pure classical states, measurement is deterministic. Probability enters classical mechanics only through mixed states (statistical ensembles, e.g., in statistical mechanics), where $\sigma_A$ will generally be a measure with positive variance, reflecting our ignorance of the precise pure state.

In quantum mechanics, by contrast, even for a \emph{pure state} (e.g., a state vector $|\psi\rangle$), the result of measuring most observables $A$ is probabilistic, and $\sigma_A$ will have positive variance (unless $|\psi\rangle$ is an eigenstate of $A$). This probabilistic nature is inherent to quantum mechanics.

\item Construction of Mixed States from Pure States: There is a map
\[ \text{Prob}(\PS_0)\longrightarrow \mcS  \]
(where $\PS_0$ is the set of pure states) that averages pure states over a probability measure to produce a general (mixed) state. If $\mu$ is a probability measure on $\PS_0$, then a state $\sigma_\mu \in \mcS$ can be defined by $\sigma_\mu = \int_{\PS_0} \pi \, d\mu(\pi)$. This expresses the idea that any state can be seen as a statistical mixture of pure states. This is formally captured by Choquet theory for compact convex sets, where states are represented by barycentric integrals over extreme points.

\item Symmetries of DOSO: A symmetry of the \dq{Data of States and Observables} (DOSO) is a collection of automorphisms of $\mcS$ and of $\mcO$ that preserve all the defining structures and compatibilities: the convex structure of $\mcS$, the algebraic and topological structure of $\mcO$, the real structure, the Lie bracket on $\mcO^\infty$, the functional calculus, and the measurement pairing. The group of all such symmetries is typically very large. Physical theories often study specific subgroups of this full symmetry group, which correspond to physical invariances (e.g., Poincaré symmetry in relativistic theories). Wigner's theorem in quantum mechanics provides a useful characterization of symmetries.
\end{enumerate}
\end{remark}

\subsubsection{Motion and Mechanical Systems}

The general DOSO framework provides the static structure. Dynamics is introduced by selecting a specific notion of \dq{motion.}

\begin{definition}[Motion and Mechanical System] \label{def:1.11}
Suppose a DOSO (Axiom System \ref{ax:DOSO}) is given.
\begin{itemize}
\item A \emph{motion} is the data of compatible one-parameter groups of automorphisms of $\mcS$ and of $\mcO$. These are denoted
\[ \sigma\mapsto\sigma^{(t)} \]
\[ A\mapsto A^{(t)} \]
for $t \in \R$.
\item The compatibility condition is that these flows are adjoint with respect to the measurement pairing:
\[ {\sigma^{(t)}}_{A}=\sigma_{A^{(t)}} \]
\item A \emph{mechanical system} is a DOSO together with a fixed, specified motion.
\end{itemize}
\end{definition}


A \dq{motion} singles out a particular way the system evolves, typically interpreted as time evolution. A mechanical system is thus not just the set of possible states and measurements, but also the rule governing how these change over time. We have ensured that this rule is physically consistent regardless of whether we view states or observables as evolving.


\begin{remark}[Properties of Motion] \label{rem:1.14}
\mbox{} % Empty box for correct spacing after remark title
\begin{enumerate}
\item Generation by a Hamiltonian: In many important cases (Hamiltonian systems), the flows defining the motion are generated by a specific real observable $H \in \mcO_{\R}^\infty$, called the \emph{Hamiltonian} of the system, via the mechanism described in Axiom 4. The Hamiltonian represents the total energy of the system in many physical contexts. If the flow is generated by $H$, then $H$ itself is conserved under this flow (i.e., $H^{(t)} = H$, or $[H,H]=0$). % The sign convention relating $H$ to the generator of the flow (e.g., whether the evolution is $e^{-itH/\hbar}$ or $e^{itH/\hbar}$) needs to be specified, we take the positive one ($e^{itH/\hbar}$).

\item Time Translation Group: The flows typically represent an action of the additive group $\R$ of real numbers, interpreted as time translations. This is the idea that the laws of physics are invariant under shifts in time.

\item Structure of Time: Time is modeled as an affine line $\mathbb{A}^1(\R)$, whose group of translations is $\R$. For simplicity, one often uses the standard real line $\R$. More sophisticated models might involve relativistic notions where time is part of a spacetime manifold, but for non-relativistic mechanics, this simple model of a uniform, universal time is usually sufficient. The affine structure means there's a notion of duration, but no preferred origin. If a Euclidean structure is added, one has a metric for time intervals.

\item \textit{Symmetries of a Mechanical System:} The symmetry group of a \emph{mechanical system} is a subgroup of the full symmetry group of its underlying DOSO. In addition to preserving the static structures, these symmetries must also preserve the specific motion (dynamics). That is, if $\alpha$ is a symmetry transformation, it must commute with the time evolution: $\alpha(\sigma^{(t)}) = (\alpha(\sigma))^{(t)}$ and $\alpha(A^{(t)}) = (\alpha(A))^{(t)}$ (or, if the symmetry itself is time-dependent, a more general condition applies). For Hamiltonian systems, this often means the symmetry transformation must commute with the Hamiltonian, or leave the Hamiltonian invariant.

The possibility of time-reversal symmetries requires special treatment (e.g., they are often realized by anti-unitary operators in quantum mechanics), which we will not do now.

Often, a particular state $\sigma \in \mcS$ (e.g., a vacuum state or ground state) might be fixed as part of the data of a mechanical system, and symmetries might be required to preserve this special state.
\end{enumerate}
\end{remark}

\begin{definition}[Stationary State] \label{def:1.15}
Fix a mechanical system with its defined motion $\sigma \mapsto \sigma^{(t)}$. A state $\sigma \in \mcS$ is called a \emph{stationary state} (or an equilibrium state, or a time-invariant state) if it does not change under this motion, i.e.,
\[ \sigma^{(t)} = \sigma \quad \text{for all } t \in \R. \]
\end{definition}

Stationary states are very important. In quantum mechanics, eigenstates of the Hamiltonian are stationary states (their corresponding density operators are invariant under unitary evolution generated by the Hamiltonian). In statistical mechanics, equilibrium ensembles (microcanonical, canonical, grand canonical) are described by stationary states. Identifying and characterizing stationary states is often a primary goal when studying a mechanical system.

\begin{remark}
    The DOSO framework, along with its compatibility conditions and the definition of a mechanical system, provides a sophisticated and abstract foundation for a vast range of physical theories. Its power lies in its generality, allowing for a unified discussion of concepts in classical, statistical, and quantum mechanics.
\end{remark}


\subsection{DOSO for Classical Mechanics}

We now illustrate how to apply the DOSO formulation to classical, non-relativistic mechanics. This serves not only to ground the axiomatic framework in familiar territory but also to highlight the unifying power of the DOSO formulation. The chosen example, a single particle in Euclidean space, while elementary, is sufficiently rich to exhibit the important features.

\subsubsection{A Particle in Euclidean Space}

For concreteness, consider a single particle of constant mass $m \in \R^{>0}$ moving in a standard $d$-dimensional Euclidean space $\mathbb{E}^d$, where $d \in \Z^{>0}$. The particle's motion is impacted by a potential energy function $V: \mathbb{E}^d \to \R$. For the purposes of defining smooth dynamics, we typically assume $V$ to be sufficiently regular, e.g., $V \in C^k(\mathbb{E}^d, \R)$ for $k \ge 2$, often $V \in C^\infty(\mathbb{E}^d, \R)$. For convenience, we denote the configuration space $M \coloneqq \mathbb{E}^d$.

A \emph{motion} or \emph{trajectory} of the particle is a sufficiently regular map $x: \R \to M$, $t \mapsto x(t)$, where $t$ represents time. The space of all such conceivable motions (prior to imposing physical laws) can be taken as $\mcF \coloneqq C^k(\R, M)$ for appropriate $k$; if we desire maximal regularity, we consider $\mcF = C^\infty(\R, M)$. This $\mcF$ can be endowed with the structure of an infinite-dimensional Fréchet manifold, but we will not say more about that here.

Physical, or \emph{classical}, trajectories are those that satisfy Newton's second law:
\[ m\ddot{x}(t) = -\nabla V(x(t)) \quad  \]
where $\ddot{x}(t) = \frac{\dd^2 x}{\dd t^2}(t)$ is the acceleration and $\nabla V(x(t))$ is the gradient of $V$ evaluated at the particle's position $x(t)$. Newton's second law is a system of $d$ second-order ordinary differential equations. An important assumption for simplifying the structure of the space of solutions is that for any given initial conditions $(x_0, v_0) \in M \times \R^d$ (initial position and initial velocity), there exists a unique solution $x(t)$ to Newton's second law for all $t \in \R$. This global existence and uniqueness is a strong condition, not always satisfied (e.g., for potentials that lead to finite-time singularities). However, when it holds, the set $N \subset \mcF$ of all physical trajectories forms a $2d$-dimensional differentiable manifold.

This manifold $N$ can be explicitly parameterized. For any chosen reference time $t_0 \in \R$, the map
\begin{align*}
\Psi_{t_0}: N &\longrightarrow M \times \R^d \cong \mathbb{E}^d \times \R^d \\
x(\cdot) &\longmapsto (x(t_0), \dot{x}(t_0))
\end{align*}
is a diffeomorphism. This map, which we may call the \emph{initial data map} at time $t_0$, provides a global chart for $N$. It's important to note that while $N$ itself is  defined as the solution space, the specific parameterization explicitly breaks the inherent time-translation symmetry of Newton's equations (if $V$ is time-independent). The space $\mathbb{E}^d \times \R^d$ can be identified with the tangent bundle $TM$ of $M=\mathbb{E}^d$.

\subsubsection{DOSO Data for the Classical Particle}

We now proceed to write the components of DOSO for this classical system.

In this deterministic classical framework, a state of maximal information corresponds to knowing the entire trajectory of the particle. Thus, the space of pure states is precisely the manifold of classical trajectories:
\[ \PS \coloneqq N \]
Each point $x(\cdot) \in N$ represents a unique, complete history of the particle satisfying the laws of motion.

For classical systems, mixed states arise from statistical uncertainty about the pure state. The encompassing space of states $\mcS$ is the space of all Borel probability measures on $N$:
\[ \mcS \coloneqq \text{Prob}(N) \]
This space $\mcS$ is a convex set. A pure state $x_N \in N$ can be identified with the Dirac delta measure $\delta_{x_N} \in \mcS$. These are the extreme points of $\mcS$.
A general mixed state $\sigma \in \mcS \setminus \PS$ describes a statistical ensemble where the particle follows one of the classical trajectories in $N$ according to the probability distribution $\sigma$. This is the domain of classical statistical mechanics.

Recall that an observable is a physical quantity whose value can be determined if the state of the system is known. In the pure state picture, an observable is a function on $N$.
\begin{itemize}
\item \textit{The Algebra $\mcO$ and its \dq{Smooth} Subspace $\mcO^\infty$:} We define the space of complex-valued observables as
\[ \mcO \coloneqq \text{Borel}(N; \C) \]
the space of complex-valued Borel-measurable functions on $N$. This is a complex vector space. The subspace of \dq{smooth} observables is
\[ \mcO^{\infty} \coloneqq C^{\infty}(N;\C) \]
the space of complex-valued smooth functions on $N$. Given that $N$ is a $2d$-dimensional smooth manifold, $C^\infty(N;\C)$ is well-defined and forms a dense subspace of $\mcO$ under suitable topologies (e.g., $L^p$-like topologies if $N$ carries a natural measure, or in the sense of approximating continuous functions). Endowing $\mcF$ (the space of all paths) with an infinite-dimensional manifold structure is a more complicated topic that we will not do. For defining observables that depend only on the physical trajectory (elements of $N$), restricting to functions on $N$ is sufficient and avoids these complexities.

\item \textit{Real Structure:} The real structure $A \mapsto A^*$ on $\mcO$ is given by pointwise complex conjugation: if $A \in \mcO$, then $A^*(x_N) \coloneqq \overline{A(x_N)}$ for all $x_N \in N$. Consequently, the space of real observables $\mcO_{\R}$ (the fixed points of $*$) is
\[ \mcO_{\R} = \text{Borel}(N;\R) \]
These are the functions that correspond to physically measurable quantities, yielding real numbers.
Similarly, $\mcO^\infty_\R = C^\infty(N;\R)$.

\item \textit{Lie Algebra Structure:} The space $\mcO^\infty_\R = C^\infty(N;\R)$ is endowed with a real Lie algebra structure via the \emph{Poisson bracket}, which requires the symplectic structure on $N$ discussed below. For $A, B \in C^\infty(N;\R)$, their Poisson bracket $\{A,B\} \in C^\infty(N;\R)$. The complexification $\mcO^\infty = C^\infty(N;\C)$ then becomes a complex Lie algebra. The compatibility $[A_1,A_2]^* = [A_1^*,A_2^*]$ is satisfied by the Poisson bracket since $\{A,B\}^* = \overline{\{A,B\}}$, and if $A^*=A, B^*=B$, then $\{A,B\}$ is real, so $\{\bar{A},\bar{B}\} = \overline{\{A,B\}}$.
\end{itemize}

The measurement pairing connects observables to probability distributions of outcomes for a given state. For $\sigma \in \mcS = \text{Prob}(N)$ and $A \in \mcO_{\R} = \text{Borel}(N;\R)$, the probability measure $\sigma_A \in \text{Prob}(\R)$ is given by the \emph{pushforward measure}:
\[ \sigma_A \coloneqq A_*\sigma\]
That is, for any Borel set $E \subset \R$, $(A_*\sigma)(E) \coloneqq \sigma(A^{-1}(E))$.
\begin{itemize}
\item \textit{Deterministic Measurement for Pure States:} If $\sigma = \delta_{x_N}$ is a pure state (corresponding to a specific trajectory $x_N \in N$), then $\sigma_A = (A_*)(\delta_{x_N}) = \delta_{A(x_N)}$. This is a point measure concentrated at the value $A(x_N)$. Thus, for pure states in classical mechanics, the measurement of any observable yields a definite, deterministic outcome. This is a key distinction from quantum mechanics.
\item \textit{Functional Calculus:} The functional calculus $f, A \mapsto f(A)$ is simply function composition: if $A \in \mcO_\R$ and $f \in \text{Borel}(\R;\R)$, then $f(A)$ is the function $x_N \mapsto f(A(x_N))$, which is also in $\mcO_\R$. The compatibility $\sigma_{f(A)} = f_* \sigma_A$ is a standard property of pushforward measures: $(f \circ A)_* \sigma = f_* (A_* \sigma)$.
\end{itemize}

The motion described in DOSO is induced by a flow on the manifold of pure states $N$. This flow, in turn, is governed by Hamiltonian mechanics, which relies on a symplectic structure on $N$.
The time evolution $x_N \mapsto (x_N)^{(t)}$ on $N$ (where $(x_N)^{(t)}(s) = x_N(s+t)$ by time translation along the trajectory) defines a one-parameter group of automorphisms of $N$.
This induces:
\begin{itemize}
\item Automorphisms on $\mcS = \text{Prob}(N)$: $\sigma \mapsto \sigma^{(t)}$, where $\sigma^{(t)}(S) = \sigma(\{x_N \mid (x_N)^{(t)} \in S\})$.
\item Automorphisms on $\mcO = \text{Borel}(N;\C)$: $A \mapsto A^{(t)}$, where $A^{(t)}(x_N) = A((x_N)^{(-t)})$.
\end{itemize}
This motion is generated by a Hamiltonian function $H \in C^\infty(N;\R)$ associated with the system's energy. The Hamiltonian flow on $N$ is a one-parameter group of symplectomorphisms. The symplectic structure itself is therefore important. We will discuss this in more detail in Lecture 2.

Without going into too much detail, a pair $(N, H)$, where $N$ is a symplectic manifold (the space of pure states) and $H \in C^\infty(N;\R)$ is a Hamiltonian function, determines a complete classical mechanical system in the DOSO framework. For our particle example, the symplectic structure on $N$ can be derived from a more foundational \emph{variational principle} (the principle of least action) applied to a Lagrangian function. This is the \emph{Lagrangian formulation}. The transition to $(N, H)$ is the \emph{Hamiltonian formulation}. The derivation of the symplectic structure within the Lagrangian framework is detailed in Lecture 3. Here, we shall merely state the symplectic 2-form on $N$.

\begin{remark}[General Symplectic Flows] \label{rem:1.20}
More generally, a classical mechanical system can be defined by a pair $(N, \{\varphi_t\}_{t\in\R})$, where $N$ is a symplectic manifold and $\{\varphi_t\}$ is a one-parameter group of symplectomorphisms (diffeomorphisms preserving the symplectic form $\omega$). Such a system always fits the DOSO framework. However, not every such flow $\{\varphi_t\}$ is \emph{Hamiltonian} in the sense of being generated by a time-independent Hamiltonian function $H \in C^\infty(N;\R)$ (i.e., $\varphi_t$ being the flow of the Hamiltonian vector field $X_H$). For instance, systems with explicitly time-dependent Hamiltonians in the traditional sense give rise to symplectic flows that are not necessarily Hamiltonian on $N$ (which is the space of solutions of an autonomous system). In Lecture 3 we will see a non-Hamiltonian symplectic flow.
\end{remark}

\subsubsection{The Symplectic Form for the Particle Example}

The symplectic structure on $N$ (the space of physical trajectories) is naturally inherited from a canonical structure on the cotangent bundle $T^*M$ of the configuration space $M = \mathbb{E}^d$.

Let $\mcF = C^\infty(\R, M)$ be the ambient space of smooth paths. One can define a presymplectic 2-form\footnote{A \textit{presymplectic 2-form} is a closed differential 2-form of constant rank on a manifold} on $\mcF$. Let $\delta$ denote the exterior derivative (variational derivative) on this infinite-dimensional manifold $\mcF$. The evaluation map is
\begin{align*}
e: \mcF \times \R &\longrightarrow M \\
(x(\cdot), t) &\longmapsto x(t)
\end{align*}
Consider variations $\delta e$ (variation of position) and $\delta \dot{e}$ (variation of velocity). A 2-form on $\mcF \times \R$ (or, by fixing time, on $\mcF$) can be defined as
\[ \Omega = m \int \langle \delta \dot{x}(t) \wedge \delta x(t) \rangle_{\R^d} \, \dd t \]
A more direct approach for the manifold of solutions $N$ is to use the parameterization $\Psi_{t_0}: N \stackrel{\cong}{\longrightarrow} TM \cong \mathbb{E}^d \times \R^d$ given by $(x(t_0), \dot{x}(t_0))$. Let $(q^1, \dots, q^d)$ be standard coordinates on $\mathbb{E}^d$ (representing $x(t_0)$) and $(v^1, \dots, v^d)$ be coordinates for the velocities $\dot{x}(t_0)$ in $\R^d$. The pullback $(\Psi_{t_0})^*\omega_{TM}$ of the canonical symplectic form $\omega_{TM}$ on $TM$ (derived from the standard Lagrangian $L(q,v) = \frac{1}{2}m\|v\|^2 - V(q)$) to $N$ is given by:
\[ \omega_N \big|_{t_0} = \sum_{i=1}^d m \, \dd v^i \wedge \dd q^i \]
This form is understood as follows: $\dd q^i$ and $\dd v^i$ are basis 1-forms on $TM$ at $(q,v)$. The important insight is that this form on $N$ is independent of the choice of $t_0$. This is a consequence of Liouville's theorem for Hamiltonian flows (which implies that the Hamiltonian flow preserves the symplectic form). Thus, $N$ carries a canonical symplectic structure.

\begin{remark}[Symplectic Form on Phase Space $T^*M$]
It is more conventional in Hamiltonian mechanics to work with phase space $T^*M$, with canonical coordinates $(q^i, p_i)$, where $p_i = \frac{\partial L}{\partial \dot{q}^i}$ are the canonical momenta. For $L = \frac{1}{2}m\sum (\dot{q}^i)^2 - V(q)$, we have $p_i = m\dot{q}^i = mv^i$. In these $(q,p)$ coordinates, the canonical symplectic 2-form on $T^*M$ is
\[ \omega_{T^*M} = \sum_{i=1}^d \dd p_i \wedge \dd q^i \]
The form from earlier becomes $m \sum \dd v^i \wedge \dd q^i$. If we set $p_i = mv^i$, then $\dd p_i = m \dd v^i$, so the forms are identical under this identification. Thus, the space of solutions $N$ is naturally identified with $T^*M$ (or $TM$), equipped with its canonical symplectic structure.
\end{remark}

\subsubsection{Dynamical Systems: Local versus Global Perspectives}

\begin{remark}[Classical Dynamics: Short Times vs. Long Times, Particles vs. Fields] \label{rem:1.24}
A symplectic manifold $(N, \omega)$ together with a smooth function $H:N\rightarrow\mathbb{R}$ (the Hamiltonian) constitutes a \emph{Hamiltonian dynamical system}. The dynamics are governed by Hamilton's equations, which can be written as $\dot{z}(t) = X_H(z(t))$, where $z(t)$ is a curve in $N$ and $X_H$ is the \emph{Hamiltonian vector field} associated with $H$, uniquely defined by the relation $i_{X_H}\omega = -\dd H$.

Local Existence and Uniqueness: The theory of ordinary differential equations (e.g., the Picard-Lindelöf theorem, applicable since $X_H$ is smooth if $H$ is $C^2$) guarantees the local existence and uniqueness of a flow $\varphi_t: N \to N$ that integrates $X_H$. This addresses the behavior of the system for \emph{small times} $t$. Constructing this local flow is relatively straightforward from a mathematical standpoint, often relying on fixed-point theorems.

Global Dynamical Questions: The more interesting and challenging questions in dynamical systems theory concern the \emph{long-time} behavior of these flows. These include:
\begin{itemize}
\item Existence of Periodic Orbits: Do trajectories exist that return to their initial state after a finite time? These are important to understanding stability and resonances.
\item Stability of Orbits: How do trajectories behave under small perturbations of initial conditions (Lyapunov stability)?
\item Ergodicity and Mixing: Does the system explore the available phase space uniformly over long times?
\item Chaotic Dynamics: Does the system exhibit sensitive dependence on initial conditions, leading to complex, unpredictable long-term behavior despite its deterministic nature? Hamiltonian systems, despite preserving phase space volume (Liouville's theorem, a consequence of $\mathcal{L}_{X_H}\omega = 0$), can indeed be chaotic (e.g., KAM theory, Smale horseshoes).
\item Limit Cycles: These are isolated periodic orbits towards which nearby trajectories converge. While prominent in dissipative systems, classical Hamiltonian systems (which are volume-preserving) do not possess attracting limit cycles in the usual sense. However, more complex recurrent structures exist.
\end{itemize}
Observe that the foundational local existence and uniqueness concerns operate at small time scales, whereas these deeper dynamical questions probe the system's behavior as $t \to \infty$.

Analogy with Classical Field Theory: Classical field theory (describing systems with infinitely many degrees of freedom, like electromagnetism or fluid dynamics) presents something similar, now involving both time and space.
\begin{itemize}
\item Local/Short-Range Questions: Here, \dq{local} means small regions in spacetime (small time differences and small spatial distances). The dynamics are typically governed by partial differential equations (PDEs), such as wave equations, Maxwell's equations, or equations from fluid dynamics (e.g., Euler or Navier-Stokes). The interesting mathematical questions about local behavior are all similar: given suitable initial data, does a unique solution exist for a short time, and does it depend continuously on the initial data? This is already a highly non-trivial area of PDE theory.
\item Global/Long-Range Questions: These pertain to the behavior of solutions for all times and over large spatial domains. Examples include: Does a solution exist globally in time, or do singularities (e.g., wave collapse, shock formation in fluids, gravitational singularities) develop in finite time? What is the asymptotic behavior of fields as $t \to \infty$? How do spatially extended structures form and evolve?
\end{itemize}

% TODO: COME BACK TO THIS
Analogy with Quantum Field Theory (QFT): The questions of QFT can often be categorized into kinematics and dynamics.
\begin{itemize}
\item Kinematics: This involves defining the important objects: quantum fields, state spaces, and observables. A key challenge at short distances (high energies/momenta) is the appearance of infinities (UV divergences) in perturbative calculations, necessitating the very complicated machinery of \emph{renormalization}. The behavior of correlation functions at short distances (operator product expansions) is important.
\item Dynamics: Much of this concerns the behavior of the theory at large distances (low energies/momenta) and long times. Key questions include: The structure of the vacuum state, the particle spectrum (existence of a mass gap), scattering theory (S-matrix), phase transitions and critical phenomena, and long-range forces and phenomena like confinement in quantum chromodynamics.

In relativistic QFTs, due to Lorentz invariance, long-time and large-distance behaviors are related. The understanding of these long-range phenomena often leads to the most interesting physical insights and drives the development of powerful non-perturbative mathematical techniques.
\end{itemize}

The long-range questions in classical mechanics, classical field theory, and quantum field theory are often the most intriguing, with significant mathematical implications. These questions frequently lead to unexpected connections between physics and seemingly unrelated areas of mathematics.

\end{remark}

\subsection{DOSO for Quantum Mechanics}

We now transition our focus to quantum mechanics, illustrating how its foundational structure can be also be encapsulated via DOSO.

A notable simplification adopted in this initial presentation is the restriction to \emph{bounded observables}. This is a common starting point in many rigorous treatments as it allows the full power of C*-algebra theory to be deployed. However, it is important to recognize that many of the most important physical observables in quantum mechanics, such as position, momentum, and typically the Hamiltonian itself (which generates time evolution), are represented by \emph{unbounded} self-adjoint operators. The theory of unbounded operators is significantly difficult. Nevertheless, the framework for bounded observables provides a good foundation and its principles often extend to the unbounded case.

\subsubsection{Hilbert Spaces and Planck's Constant}

At the center of quantum theory lies a universal constant of nature, \textit{Planck's constant} $\hbar$. This constant has the physical dimensions of \emph{action}:
\[ [\hbar] = \frac{\text{mass} \cdot \text{length}^2}{\text{time}}\]
Recalling that energy has units of $\text{mass} \cdot \text{length}^2/\text{time}^2$, the units of $\hbar$ are equivalently those of energy $\cdot$ time. Planck's constant quantifies the important scale at which quantum effects become manifest, appearing often in commutation relations (e.g., $[x,p]=i\hbar$), uncertainty principles (e.g., $\Delta x \Delta p \ge \hbar/2$), and the quantization of physical quantities (e.g., energy levels of bound systems, angular momentum).

The entire DOSO structure of a quantum mechanical system is derived from a single, foundational mathematical object: The Hilbert space. To every quantum mechanical system is associated a \textit{separable complex Hilbert space} $\mcH$, which may be finite or countably infinite dimensional. This space $\mcH$ is the place where the states and observables of the system are defined.

Complex numbers are essential in quantum mechanics for describing superposition, interference, and time evolution via the Schrödinger equation. Probability amplitudes, central to the Born rule, rely on the complex structure of Hilbert space. The Hilbert space's inner product defines measurement probabilities, orthogonality, and self-adjoint operators (observables), while completeness ensures convergence and enables the spectral theorem. Foundational quantum systems typically use separable Hilbert spaces with countable orthonormal bases, though non-separable spaces appear in quantum field theory.

\subsubsection{DOSO Data for Quantum Systems}

With the Hilbert space $\mcH$ established, we can now define the components of the DOSO framework for a quantum system.

In quantum mechanics, a pure state corresponds to a \emph{ray} in the Hilbert space $\mcH$.
\begin{definition}
     A \textit{ray} is an equivalence class of non-zero vectors differing by a non-zero complex scalar multiple: $\tilde{\psi} = \{c\psi \mid c \in \C, c \neq 0\}$, where $\psi \in \mcH \setminus \{0\}$. 
\end{definition}
Conventionally, rays are represented by unit vectors $\|\psi\|=1$, with the understanding that vectors differing only by a phase factor $e^{i\alpha}$ (an unobservable global phase) represent the same physical state.
The space of all such rays is the \textit{projective Hilbert space} $P(\mcH)$.
\[ \PS \coloneqq P(\mcH) \]

The general space of states $\mcS$, encompassing both pure and mixed states, is the set of \textit{density operators} on $\mcH$. 
\begin{definition}
    A \textit{density operator} $\rho$ is a linear operator on $\mcH$ satisfying:
\begin{enumerate}
\item $\rho$ is self-adjoint: $\rho = \rho^*$.
\item $\rho$ is non-negative (positive semi-definite): $\langle\phi, \rho\phi\rangle \ge 0$ for all $\phi \in \mcH$.
\item $\rho$ is trace-class, and its trace is unity: $\Tr(\rho) = 1$.
\end{enumerate}
\[ \mcS \coloneqq \{ \rho \in \mathcal{T}_1(\mcH) \mid \rho = \rho^*, \rho \ge 0, \Tr(\rho)=1 \} \]
where $\mathcal{T}_1(\mcH)$ is the space of trace-class operators on $\mcH$.
\end{definition}
The set $\mcS$ is a convex subset of the real Banach space of self-adjoint trace-class operators. A density operator $\rho$ represents a pure state if and only if $\rho^2 = \rho$ (in addition to the other defining properties), or equivalently, if $\Tr(\rho^2) = 1$. Otherwise, if $\Tr(\rho^2) < 1$, $\rho$ represents a mixed state, which can be expressed as a convex combination of pure states (e.g., $\rho = \sum_k p_k \ket{\psi_k}\bra{\psi_k}$ with $p_k > 0, \sum p_k = 1$).

\begin{itemize}
\item \textit{The Algebra $\mcO$ and $\mcO^\infty$: Bounded Operators on $\mcH$}\\
In this simplified exposition focusing on bounded observables, we define:
\[ \mcO = \mcO^\infty \coloneqq \text{End}(\mcH) \equiv \mathcal{B}(\mathcal{H}) \]
the complex Banach space (in fact, a C*-algebra) of all bounded linear operators mapping $\mcH$ to itself. The restriction to bounded operators ensures that they are defined on the entirety of $\mcH$ and possess desirable analytic properties.

\item \textit{Real Structure: Operator Adjoint}\\
The real structure $A \mapsto A^*$ on $\mcO = \mathcal{B}(\mathcal{H})$ is given by the standard Hilbert space \textit{adjoint operation}. An operator $A^*$ is the adjoint of $A$ if $\langle\phi, A\psi\rangle = \langle A^*\phi, \psi\rangle$ for all $\phi, \psi \in \mcH$. This map is antilinear and involutive.

\item \textit{Real Observables $\mcO_\R$: Self-Adjoint Operators}\\
Consequently, the space of real observables $\mcO_\R$ consists of all bounded \textit{self-adjoint operators} on $\mcH$:
\[ \mcO_\R = \{ A \in \mathcal{B}(\mathcal{H}) \mid A = A^* \} \]
These operators have real spectrum and correspond to physically measurable quantities whose outcomes are real numbers.

\item \textit{Lie Algebra Structure: The Quantum Commutator}\\
The complex Lie algebra structure on $\mcO = \mathcal{B}(\mathcal{H})$ is defined via a scaled commutator:
\[ [A_1, A_2]_\hbar \coloneqq \frac{-i}{\hbar} (A_1 A_2 - A_2 A_1)  \]
for $A_1, A_2 \in \mathcal{B}(\mathcal{H})$. Here, $A_1 A_2$ denotes operator composition.

If $A_1$ and $A_2$ are self-adjoint, their commutator $C = A_1 A_2 - A_2 A_1$ is skew-adjoint ($C^* = -C$). Multiplication by $-i/\hbar$ (or $1/(i\hbar)$) then makes $[A_1, A_2]_\hbar$ self-adjoint. Thus, this definition ensures that the Lie bracket of two real (self-adjoint) observables is itself a real (self-adjoint) observable, making $\mcO_\R$ a real Lie algebra under this bracket. This Lie algebra structure is the quantum analogue of the Poisson bracket algebra in classical mechanics and is important to the Heisenberg uncertainty principle and to the equations of motion. Thus, the compatibility condition $[X,Y]^* = [X^*,Y^*]$ is satisfied by this definition.


\item \textit{Functional Calculus for Observables}\\
Functions of observables are defined via the \textit{spectral theorem} for self-adjoint operators, which we discuss in Lecture 4. If $A \in \mcO_\R$ is a bounded self-adjoint operator and $f: \R \to \R$ is a Borel function, then $f(A)$ is a well-defined bounded self-adjoint operator. This is important for constructing observables like $A^2$ or $\exp(A)$, and particularly for defining projection-valued measures $\Pi_A(E) \coloneqq \chi_E(A)$ (where $\chi_E$ is the characteristic function of a Borel set $E \subset \R$), which are central to the theory of measurement.
\end{itemize}

The measurement pairing connects an observable $A \in \mcO_\R$ and a state (pure or mixed) to a probability distribution $\sigma_A \in \text{Prob}(\R)$ for the outcomes of measuring $A$. Let $\Pi_A$ denote the projection-valued measure associated with the self-adjoint operator $A$ via the spectral theorem (so $\Pi_A(E)$ is the projection onto the subspace of states for which the value of $A$ lies in the Borel set $E \subset \R$).

\begin{itemize}
\item \textit{Measurement in a Pure State:}
Suppose the system is in a pure state represented by a unit vector $\psi \in \mcH$ (equivalently, by the line $L = \text{span}\{\psi\} \in P(\mcH)$ or the density operator $\rho_\psi = \ket{\psi}\bra{\psi}$). The probability measure $\sigma_A \equiv (\rho_\psi)_A$ on $\R$ is given by the \textit{Born rule}:
\[ (\rho_\psi)_A(E) \coloneqq \langle\psi, \Pi_A(E)\psi\rangle = \Tr(\rho_\psi \Pi_A(E))\]
for any Borel subset $E \subset \R$. This value is the probability that a measurement of $A$ yields a result in $E$. The independence of this probability from the choice of unit vector representing the ray $L$ is ensured by the quadratic dependence on $\psi$ (an overall phase $e^{i\alpha}$ cancels out: $\langle e^{i\alpha}\psi, \Pi_A(E)e^{i\alpha}\psi\rangle = \langle\psi, \Pi_A(E)\psi\rangle$).

\item \textit{Measurement in a Mixed State:}
If the system is in a general (possibly mixed) state described by a density operator $S \in \mcS$, then we can generalize to:
\[ S_A(E) \coloneqq \Tr(S \Pi_A(E)) \quad \]
for any Borel subset $E \subset \R$. This formula is linear in $S$, reflecting the convex nature of states, and provides the statistical predictions of quantum theory for any preparation $S$ and any (bounded) observable $A$. If $S = \sum_k p_k \ket{\psi_k}\bra{\psi_k}$, then $S_A(E) = \sum_k p_k \langle\psi_k, \Pi_A(E)\psi_k\rangle$.
\end{itemize}

\textit{Quantum Motion (Dynamics)}
The motion (time evolution) of a quantum system is generated by a specific self-adjoint observable $H \in \mcO_\R^\infty$ (here, $\mcO_\R^\infty = \mcO_\R$ are bounded self-adjoint operators), designated as the \textit{Hamiltonian} of the system. In most realistic systems, $H$ is unbounded, but the principles extend via Stone's theorem.

\begin{itemize}
\item \textit{Unitary Time Evolution Group:}
The Hamiltonian $H$ generates a one-parameter group of \textit{unitary transformations} $\{U_t\}_{t\in\R}$ on $\mcH$:
\[ U_t \coloneqq e^{-itH/\hbar}, \quad t \in \R \]
This is defined using the functional calculus for self-adjoint operators (applying $f_t(\lambda) = e^{-it\lambda/\hbar}$ to $H$). Since $H$ is self-adjoint, $U_t$ is unitary for all $t \in \R$ ($U_t^* = U_t^{-1} = U_{-t}$). These operators describe the evolution of the system over a time interval $t$. For bounded $H$, this group is norm-continuous.

\item \textit{Evolution of States (Schrödinger Picture):}
If a state at time $t=0$ is $S \equiv S(0)$, its evolution to time $t$ is given by:
\[ S^{(t)} \equiv S(t) \coloneqq U_t S(0) U_t^* \]
This preserves the properties of a density operator (self-adjointness, non-negativity, trace one). If $S(0) = \ket{\psi_0}\bra{\psi_0}$ is a pure state, then $S(t) = \ket{\psi(t)}\bra{\psi(t)}$ where $\ket{\psi(t)} = U_t\ket{\psi_0}$. The latter satisfies the \textit{Schrödinger equation}: $i\hbar \frac{\dd}{\dd t}\ket{\psi(t)} = H\ket{\psi(t)}$. More generally, density operators evolve according to the \textit{Liouville-von Neumann equation}: $i\hbar \frac{\dd}{\dd t}S(t) = [H, S(t)] \equiv HS(t) - S(t)H$.

\item \textit{Evolution of Observables (Heisenberg Picture):}
Alternatively, one can consider the states fixed and the observables evolving. An observable $A \equiv A(0)$ evolves to $A^{(t)} \equiv A(t)$ according to:
\[ A^{(t)} \equiv A(t) \coloneqq U_t^* A(0) U_t \]
These evolved observables $A(t)$ satisfy the \textit{Heisenberg equation of motion}:
\[ \frac{\dd}{\dd t}A(t) = \frac{i}{\hbar}[H, A(t)] \]
\end{itemize}
The compatibility condition $(S^{(t)})_A = S_{(A^{(t)})}$ is satisfied:
\[\Tr(S^{(t)}\Pi_A(E)) = \Tr(U_t S U_t^* \Pi_A(E)) = \Tr(S U_t^* \Pi_A(E) U_t).\] 
The term $U_t^* \Pi_A(E) U_t$ is the PVM for the evolved observable $A^{(t)}$, evaluated on $E$. That is, 
\[\Pi_{A^{(t)}}(E) = U_t^* \Pi_A(E) U_t.\] Thus 
\[\Tr(S \Pi_{A^{(t)}}(E)) = S_{(A^{(t)})}(E).\]

We have finished explaining how to axiomatize quantum mechanics using DOSO, and this is a good place to stop.